\chapter{Valutazione retrospettiva}
\label{chap:conclusioni}

\textit{Il presente capitolo propone una valutazione critica del percorso di tirocinio svolto. Attraverso l'analisi dei risultati conseguiti rispetto agli obiettivi prefissati e una riflessione sulla maturazione professionale acquisita, si intende fornire una sintesi del valore dell'esperienza in Ergon Informatica. Il capitolo si conclude con una disamina sul rapporto tra il bagaglio formativo universitario e le competenze richieste dal mondo del lavoro, tracciando un bilancio finale sul percorso accademico intrapreso.}

\section{Analisi dei risultati e maturazione professionale}
Questa sezione valuta, su base oggettiva, il grado di soddisfacimento degli obiettivi prefissati in sede di pianificazione. Si analizzano le performance ottenute confrontandole con le aspettative aziendali, evidenziando come la gestione autonoma di un progetto \textit{enterprise} — unita alla necessità di apprendere rapidamente \textit{framework} e tecnologie inedite — abbia favorito una crescita professionale significativa. L'attenzione si focalizza sulla capacità di allineare il proprio operato agli standard e alla visione dell'azienda, superando la prospettiva puramente individuale in favore di una soluzione tecnica coerente con le necessità produttive e di mercato.

\section{Bilancio formativo: dall'università al mondo del lavoro}
Viene analizzato il divario tra le competenze acquisite durante il corso di studi e quelle richieste per l'operatività quotidiana. Si discute come l'approccio ingegneristico fornito dall'Università di Padova abbia costituito la base metodologica necessaria per affrontare il tirocinio, identificando al contempo le aree di competenza che hanno richiesto un maggior adattamento. Il capitolo si chiude con una riflessione retrospettiva sul percorso accademico, riconoscendo il ruolo dei singoli insegnamenti nel fornire gli strumenti teorici necessari a comprendere, a posteriori, l'importanza della rigorosità scientifica e della pianificazione nel contesto professionale.

\newpage